\documentclass{article}
\usepackage{graphicx} % Required for inserting images
\usepackage{algorithm} %
\usepackage{algpseudocode}
\usepackage{amsfonts}

\usepackage{tcolorbox}
\newtcolorbox{greenbox}[1]{colback=green!5!white,colframe=green!75!black,fonttitle=\bfseries,title=#1}


\usepackage{amsthm}
\usepackage{bm}

\newtheorem{definition}{Definition}
\newtheorem{theorem}{Theorem}
\newtheorem{lemma}{Lemma}

\newtheorem*{remark}{Remark}

\usepackage{amsmath}
\DeclareMathOperator*{\argmin}{argmin} % no space, limits underneath in displays



\setlength{\parindent}{0pt} % for getting rid of first line indentation


% ================================================
% =====================  BEGIN DOCUMENT 
% ================================================

\title{Corporate Finance and Investment Science}
\author{Juan Pablo Bello Gonzalez}
\date{January 2024}

\begin{document}



\maketitle

\section{Financial Markets and NPV}
Financial markets exist so companies and individuals can adjust consumption across time periods. The purpose is also to facilitate borrowing and lending.

Investors can invest directly into fiems or via third parties called financial intermediaries which grant them flexibility. 

financial intermediation can take three forms: 
\begin{enumerate}
\item Size intermediation. Eg. banks combine investments of many small investors into large loans for a few borrowers. 
\item Term intermediation. take from short term lenders (savings account) and re-lend to long term borrowers (mortgages).
\item Risk Intermediation. Intermediaries can assume the excess risk and protect borrower (hedge-funds investing into crypto)
\end{enumerate}

\begin{definition}
Project. Sequence of cashflows spread over time. 
\end{definition}

To evaluate a project we can use several rules. 
\subsection{NPV}
\begin{theorem}
\textbf{NPV rule:} project is good if NPV > 0. This implies the project is better than risk-free borrowing and lending in the market. 
\end{theorem}
Pros:
\begin{itemize}
\item Takes TVM into consideration, discounts all cashflows. 
\item Easy to evaluate NPV for combinations of projects. 
\item Can be applied to any weird pattern of cashflows. 
\end{itemize}

Cons: 
\begin{itemize}
\item Cashflows must be forecasted. 
\item What should the discount rate be?
\item Doesn't take into consideration size of project. 
\item Doesn't take into consideration required projects, or projects with a differnt time horizon, like a choice between two required insurance plans. 
\end{itemize}

\begin{definition}
\textbf{Capital Budgeting} is the decision-making process for accepting or rejecting new projects/investments.
\end{definition}

\begin{theorem}
\[ \text{NPV} > 0 \iff \text{increase in company value} \iff \text{increase in shareholder value}\]
\end{theorem}

\subsection{Payback Period Rule}
Measures how long it takes to "payback" the initial investment of a project.

Pros:
\begin{itemize}
\item easy
\item gives you a measure of liquidity. 
\end{itemize}

Cons:
\begin{itemize}
\item Ignores TVM
\item Ignores all cashflows after the payback period
\item What is the comparison rule? is a payback period of 4 years good? ambiguity.
\end{itemize}

so in terms of comparing projects, this only does one thing right: give a measure of liquidity. 
\subsection{Discounted payback period rule}
Number of periods in the future when the cummulative discounted cash flows total the initial investment. 
Pros:
\begin{itemize}
\item gives a measure of liquidity. 
\item Doesn't ignore TVM
\end{itemize}

Cons:
\begin{itemize}
\item Ignores all cashflows after the payback period
\item Ambiguity of comparison rule
\end{itemize}

\subsection{Internal Rate of Return}
\begin{definition}
\textbf{Cost of capital:} minimum RATE OF RETURN ("cost") needed to justify the "capital" needed to undertake a project.
\end{definition}

\begin{definition}
Cost of capital such that the NPV of a project is 0
\end{definition}

Rule: Accept a project if IRR > required rate of return. This "required rate of return" 

\subsection{Profitability Index}






\end{document}
